% Shared: relation to prior work (what is classical, what is new).
\paragraph{Relation to prior work}
Most methods that learn a PDE from data regress a time derivative on a library of candidate terms: SINDy and
PDE-FIND~\cite{BruProKut16,RudBruProKut17}, their weak, ensemble and Bayesian variants~\cite{MesBor21,FasKutBru22,HirBarKut22},
and Gaussian-process formulations of linear operators~\cite{RaiPerKar17}. Their unit is the \emph{term}. Dynamic mode
decomposition~\cite{Sch10,SasKut22} works with eigenvalues, as we do, but returns a linear operator on the data, not an
equation, and does not group eigenvalues by the factor they belong to. Exponential fitting by matrix
pencils~\cite{HuaSar90} is a classical tool of signal processing; the algebra of constant-coefficient operators, the
Ehrenpreis--Palamodov fundamental principle and the solvability results behind \cref{thm:factor} are classical
analysis~\cite{Ehr70,Pal70,Mal56,Hor90,Obe90}; Dirichlet-process mixtures are standard Bayesian
clustering~\cite{BleJor06}. What is new here is the connection between them. (i) We identify the \emph{factor} of the
symbol, not the term, as the unit a learning method should group by, and we prove that, for two terms whose varieties meet in finitely many points, per-term grouping is impossible
outside a finite-dimensional set of solutions (\cref{prop:terms,prop:discont}). (ii) We state what finite data can and
cannot certify about such groups (\cref{thm:ident,prop:certify}) and give the exact conditioning of the
separation between groups (\cref{thm:hm,prop:bayes-angle}). (iii) We build an algorithm that groups twice, modes into
factors within a sample and factors across posterior samples, and whose output carries these limits as labels and
intervals.
